% ===========================================================================
%  Devoir surveillé nº 2 — fin de la partie Analyse
% ===========================================================================
\chapter*{Devoir surveillé nº\,2 — Analyse}
\markboth{Devoir surveillé nº\,2}{Devoir surveillé nº\,2}
\addcontentsline{toc}{chapter}{Devoir surveillé nº\,2 — Analyse}

\begin{devoirsurveille}[portée : chapitres 1 à 9]
\textbf{Durée : 2 h 15.} Les deux exercices et le problème sont obligatoires.
Calculatrice scientifique non programmable autorisée. Barème
\textup{(A1 ; A2)}.

\vspace{0.7em}
\noindent\textbf{\sffamily EXERCICE 1} \hfill \textup{(05 points)}

\begin{enumerate}[label=\textbf{\arabic*.}, leftmargin=*, itemsep=0.25em]
  \item Résoudre dans $\R$ l'équation $\ln(x+2)+\ln(x-1) = \ln 4$.
        \bareme{1,5 ; 1,5}
  \item Résoudre dans $\R$ l'équation $e^{2x}-6e^{x}+5 = 0$.
        \bareme{1,5 ; 1,5}
  \item Résoudre dans $\R$ l'inéquation $\ln(3-x) \leq \ln(x+1)$.
        \bareme{1 ; 1}
  \item Simplifier $A = \ln 8-3\ln 2+\ln\left(e^{2}\right)$.
        \bareme{1 ; 1}
\end{enumerate}

\vspace{0.7em}
\noindent\textbf{\sffamily EXERCICE 2} \hfill \textup{(05 points)}

La courbe $(\Gamma)$ ci-dessous est celle d'une fonction $g$ définie sur
$\intervalle{-4}{4}$. La droite en pointillé a pour équation $y = 2$.

\begin{center}
\begin{tikzpicture}
\begin{axis}[
  width = 10.4cm, height = 5.8cm,
  axis lines = middle, xlabel = {$x$}, ylabel = {$y$},
  xmin = -4.6, xmax = 4.8, ymin = -3.4, ymax = 4.6,
  xtick = {-4,-3,-2,-1,1,2,3,4}, ytick = {-3,-2,-1,1,2,3,4},
  axis line style = {bmtGris},
  tick label style = {font=\scriptsize, color=bmtGris},
  grid = both, grid style = {bmtGrisClair, line width=0.3pt}, clip = true,
]
  \addplot[bmtIndigo, line width=1.4pt, domain=-4:4, samples=200]
    {0.22*x^3-0.7*x^2-0.4*x+3};
  \addplot[bmtLaterite, line width=0.9pt, dash pattern=on 3pt off 2pt,
           domain=-4.5:4.7] {2};
\end{axis}
\end{tikzpicture}
\end{center}

\begin{enumerate}[label=\textbf{\arabic*.}, leftmargin=*, itemsep=0.25em]
  \item Lire $g(0)$ et $g(2)$. \bareme{0,5 ; 0,5}
  \item Dresser le tableau de variation de $g$ sur $\intervalle{-4}{4}$.
        \bareme{1,5 ; 1,5}
  \item Donner le maximum local et le minimum local de $g$, en précisant les
        abscisses où ils sont atteints. \bareme{1 ; 1}
  \item Déterminer graphiquement le nombre de solutions de l'équation
        $g(x) = 2$. \bareme{1 ; 1}
  \item Résoudre graphiquement l'inéquation $g(x) \geq 2$. \bareme{1 ; 1}
\end{enumerate}

\vspace{0.7em}
\noindent\textbf{\sffamily PROBLÈME} \hfill \textup{(10 points)}

On considère la fonction numérique $f$ définie sur
$\intervalleo{0}{+\infty}$ par
\[
  f(x) = 2x-3+2\ln x ,
\]
et l'on note $(\mathcal C)$ sa courbe représentative dans un repère
orthonormé $\left(O\,;\,\veci\,,\,\vecj\right)$ d'unité graphique $1$ cm.

\begin{enumerate}[label=\textbf{\arabic*.}, leftmargin=*, itemsep=0.25em]
  \item Calculer $\displaystyle\lim_{x \to 0^{+}} f(x)$ et en déduire une
        asymptote à $(\mathcal C)$. \bareme{1 ; 1}
  \item \begin{enumerate}[label=\textbf{\alph*)}, leftmargin=1.4em, itemsep=0.15em]
      \item Vérifier que, pour tout $x>0$,
            $f(x) = x\left(2-\dfrac3x+2\,\dfrac{\ln x}{x}\right)$.
            \bareme{1 ; 0,75}
      \item Sachant que $\limpinf \dfrac{\ln x}{x} = 0$, calculer
            $\limpinf f(x)$. \bareme{1 ; 0,75}
    \end{enumerate}
  \item \begin{enumerate}[label=\textbf{\alph*)}, leftmargin=1.4em, itemsep=0.15em]
      \item Calculer $f'(x)$ pour tout $x>0$. \bareme{1 ; 1}
      \item Étudier son signe et dresser le tableau de variation de $f$.
            \bareme{1,5 ; 1,5}
    \end{enumerate}
  \item Montrer que l'équation $f(x) = 0$ admet une solution comprise entre
        $1$ et $2$, puis en donner une valeur approchée à $0{,}1$ près à
        l'aide de la calculatrice. \bareme{1,5 ; 1,25}
  \item Écrire une équation de la tangente $(T)$ à $(\mathcal C)$ au point
        d'abscisse $1$. \bareme{1 ; 0,75}
  \item Calculer $f(0{,}5)$, $f(2)$ et $f(4)$ à $10^{-2}$ près, puis tracer
        $(T)$ et $(\mathcal C)$ sur $\intervalleof{0}{4}$. \bareme{2 ; 1,5}
\end{enumerate}

\pourAdeux
\begin{enumerate}[label=\textbf{\arabic*.}, leftmargin=*, itemsep=0.25em, start=7]
  \item Soit $F$ la fonction définie sur $\intervalleo{0}{+\infty}$ par
        $F(x) = x^{2}-5x+2x\ln x$.
    \begin{enumerate}[label=\textbf{\alph*)}, leftmargin=1.4em, itemsep=0.15em]
      \item Montrer que $F$ est une primitive de $f$ sur
            $\intervalleo{0}{+\infty}$. \bareme{ ; 0,75}
      \item Calculer $\displaystyle\int_{1}^{e} f(x)\dx$. On donne
            $e \approx 2{,}7$ et $e^{2} \approx 7{,}4$. \bareme{ ; 1}
    \end{enumerate}
\end{enumerate}
\end{devoirsurveille}

\begin{corrige}[devoir surveillé nº 2 — exercice 1]
\textbf{1.} Conditions d'existence : $x+2>0$ et $x-1>0$, donc $x>1$.
L'équation devient $\ln\left[(x+2)(x-1)\right] = \ln4$, soit
$x^{2}+x-2 = 4$, c'est-à-dire $x^{2}+x-6 = 0$. Le discriminant vaut $25$ et
les racines sont $-3$ et $2$. On rejette $-3$, qui ne vérifie pas les
conditions : $\mathcal S = \{2\}$.

\textbf{2.} En posant $X = e^{x}>0$ : $X^{2}-6X+5 = 0$, de racines $1$ et
$5$, toutes deux positives. Donc $e^{x} = 1$ ou $e^{x} = 5$, soit
$\mathcal S = \{0\,;\,\ln5\}$.

\textbf{3.} Conditions : $3-x>0$ et $x+1>0$, donc $-1<x<3$. La fonction
$\ln$ étant croissante, l'inéquation équivaut à $3-x \leq x+1$, soit
$x \geq 1$. D'où $\mathcal S = \intervallefo{1}{3}$.

\textbf{4.} $\ln8 = 3\ln2$, donc $A = 3\ln2-3\ln2+2 = 2$.
\end{corrige}

\begin{corrige}[devoir surveillé nº 2 — exercice 2]
\textbf{1.} $g(0) = 3$ et $g(2) \approx 0{,}4$.

\textbf{2.} La courbe monte de $-4$ jusqu'à environ $-0{,}25$, descend
jusqu'à environ $2{,}4$, puis remonte jusqu'à $4$ :

\begin{center}\footnotesize
\begin{tabular}{|c|ccccccc|}
\hline
$x$ & $-4$ & & $-0{,}25$ & & $2{,}4$ & & $4$\\
\hline
$g$ & $-20$ & $\nearrow$ & $3{,}05$ & $\searrow$ & $0{,}33$ & $\nearrow$ &
$3{,}48$\\
\hline
\end{tabular}
\end{center}

\textbf{3.} Maximum local $\approx 3{,}05$ atteint en $x \approx -0{,}25$ ;
minimum local $\approx 0{,}33$ atteint en $x \approx 2{,}4$.

\textbf{4.} La droite $y = 2$ coupe la courbe en trois points : l'équation
$g(x) = 2$ a trois solutions.

\textbf{5.} $g(x) \geq 2$ là où la courbe est au-dessus de la droite :
approximativement sur $\intervalle{-1{,}3}{0{,}8}$ et sur
$\intervalle{3{,}4}{4}$.
\end{corrige}

\begin{corrige}[devoir surveillé nº 2 — problème]
\textbf{1.} En $0^{+}$, $2x-3 \to -3$ et $2\ln x \to -\infty$ : donc
$f(x) \to -\infty$. L'axe des ordonnées, d'équation $x = 0$, est asymptote
verticale à $(\mathcal C)$.

\textbf{2. a)} $x\left(2-\frac3x+2\frac{\ln x}{x}\right) = 2x-3+2\ln x$.
\checkmark
\textbf{b)} La parenthèse tend vers $2-0+0 = 2$ et $x \to +\infty$ : donc
$f(x) \to +\infty$.

\textbf{3. a)} $f'(x) = 2+\dfrac2x = \dfrac{2x+2}{x} = \dfrac{2(x+1)}{x}$.
\textbf{b)} Sur $\intervalleo{0}{+\infty}$, $x>0$ et $x+1>0$ : la dérivée
est strictement positive. La fonction est donc strictement croissante sur
$\intervalleo{0}{+\infty}$, de $-\infty$ à $+\infty$.

\textbf{4.} $f$ est continue et strictement croissante ; $f(1) = -1<0$ et
$f(2) = 1+2\ln2 \approx 2{,}39>0$. Elle s'annule donc une fois et une seule
entre $1$ et $2$. La calculatrice donne $f(1{,}3) \approx -0{,}12$ et
$f(1{,}4) \approx 0{,}47$ : la solution vaut environ $1{,}3$.

\textbf{5.} $f(1) = -1$ et $f'(1) = 4$ : $(T) : y = 4(x-1)-1 = 4x-5$.

\textbf{6.} $f(0{,}5) = -2-2\ln2 \approx -3{,}39$ ;
$f(2) \approx 2{,}39$ ; $f(4) = 5+2\ln4 \approx 7{,}77$. On trace l'axe des
ordonnées comme asymptote, la tangente $y = 4x-5$, puis la courbe, qui monte
sans interruption en coupant l'axe des abscisses vers $1{,}3$.

\textbf{7. a)} $F'(x) = 2x-5+2\left(\ln x+1\right) = 2x-3+2\ln x = f(x)$.
\checkmark
\textbf{b)} $\displaystyle\int_{1}^{e} f(x)\dx = F(e)-F(1)
= \left(e^{2}-5e+2e\right)-\left(1-5+0\right) = e^{2}-3e+4
\approx 7{,}4-8{,}1+4 = 3{,}3$.
\end{corrige}

\begin{notepourlaclasse}
Le problème est construit sur le modèle exact des sujets récents, à une
différence près : la fonction y est strictement monotone, ce qui simplifie
le tableau de variation et libère du temps pour la question 4, celle de
l'encadrement de la racine. Cette question n'apparaît pas dans tous les
sujets malgaches, mais elle est la plus formatrice du devoir.

Si le temps manque, supprimez la question 6 (tracé) et ramenez le devoir à
une heure trente : il reste seize points et toutes les techniques sont
couvertes.

L'exercice 2 est volontairement entièrement graphique. Les élèves qui
échouent sur les calculs d'analyse y retrouvent quatre points accessibles —
c'est un levier de motivation à ne pas négliger en cours d'année.
\end{notepourlaclasse}
